\subsubsection*{\texttt{25-index-sentwindow-markdown}}
\label{app:samples:25_index_sentwindow_markdown}

\paragraph{Q3: How long does the Haste spell last?}
\noindent \textbf{Corr.:}~5.00/5\quad \textbf{Faith.:}~1.00\quad \textbf{Relev.:}~1.00\quad \textbf{Ctx.Rel.:}~1.00\\

\noindent\textbf{Antwort:} The Haste spell lasts for a duration of concentration, up to 1 minute.

\noindent\textbf{Abgerufene Chunks:}
\begin{itemize}[noitemsep, topsep=2pt]
  \item[\textbf{\#1}] \texttt{[0.0167]} \textit{Spells/Modify Memory} -- While this charm lasts, you can affect the target's memory of an event that it e\ldots
  \item[\textbf{\#2}] \texttt{[0.0164]} \textit{Spells/Haste} -- \#\#\# Haste  *3rd-level transmutation*  **Casting Time:** 1 action  **Range:** 30 \ldots
  \item[\textbf{\#3}] \texttt{[0.0161]} \textit{Spells/Bestow Curse} -- If you use a spell slot of 5th level or higher, the duration is 8 hours.
  \item[\textbf{\#4}] \texttt{[0.0159]} \textit{Spells/Bestow Curse} -- If you use a spell slot of 7th level or higher, the duration is 24 hours.
  \item[\textbf{\#5}] \texttt{[0.0156]} \textit{Spells/Seeming} -- The spell lasts for the duration, unless you use your action to dismiss it soone\ldots
\end{itemize}
\vspace{4pt}
\hrule
\vspace{6pt}

\paragraph{Q5: Can a grappled creature cast spells?}
\noindent \textbf{Corr.:}~1.00/5\quad \textbf{Faith.:}~0.00\quad \textbf{Relev.:}~0.00\quad \textbf{Ctx.Rel.:}~0.25\\

\noindent\textbf{Antwort:} I cannot find this in the provided rules about whether a grappled creature can cast spells or not.

\noindent\textbf{Abgerufene Chunks:}
\begin{itemize}[noitemsep, topsep=2pt]
  \item[\textbf{\#1}] \texttt{[0.0167]} \textit{Characterizations/Feats} -- - You can use your action to try to pin a creature grappled by you.
  \item[\textbf{\#2}] \texttt{[0.0164]} \textit{Gameplay/Combat} -- \#\# Cast a Spell  Spellcasters such as wizards and clerics, as well as many monst\ldots
  \item[\textbf{\#3}] \texttt{[0.0161]} \textit{Spells/# Spellcasting} -- They might turn an enemy into a harmless creature, bolster the strength of an al\ldots
  \item[\textbf{\#4}] \texttt{[0.0159]} \textit{Spells/# Spellcasting} -- Such spells can make enemies see the caster as a friend, force creatures to take\ldots
  \item[\textbf{\#5}] \texttt{[0.0156]} \textit{Monsters/Couatl} -- It can innately cast the following spells, requiring only verbal Components  At \ldots
\end{itemize}
\vspace{4pt}
\hrule
\vspace{6pt}

\paragraph{Q12: If a creature has exhaustion level 3 and is also restrained, what is its movement speed?}
\noindent \textbf{Corr.:}~5.00/5\quad \textbf{Faith.:}~1.00\quad \textbf{Relev.:}~1.00\quad \textbf{Ctx.Rel.:}~0.25\\

\noindent\textbf{Antwort:} The creature's movement speed is 0 because of the restrained condition.

\noindent\textbf{Abgerufene Chunks:}
\begin{itemize}[noitemsep, topsep=2pt]
  \item[\textbf{\#1}] \texttt{[0.0167]} \textit{Gamemastering/Conditions} -- \#\# Restrained  - A restrained creature's speed becomes 0, and it can't benefit f\ldots
  \item[\textbf{\#2}] \texttt{[0.0164]} \textit{Gameplay/Adventuring} -- For example, a creature with a Constitution of 14 can hold its breath for 3 minu\ldots
  \item[\textbf{\#3}] \texttt{[0.0161]} \textit{Gameplay/Adventuring} -- \#\# Suffocating  A creature can hold its breath for a number of minutes equal to \ldots
  \item[\textbf{\#4}] \texttt{[0.0159]} \textit{Monsters/Behir} -- If the behir dies, a swallowed creature is no longer restrained by it and can es\ldots
  \item[\textbf{\#5}] \texttt{[0.0156]} \textit{Monsters/Kraken} -- It can spend 5 feet of movement to escape from nonmagical restraints or being gr\ldots
\end{itemize}
\vspace{4pt}
\hrule
\vspace{6pt}

\paragraph{Q16: Can you use Action Surge to cast two leveled spells in one turn?}
\noindent \textbf{Corr.:}~1.00/5\quad \textbf{Faith.:}~1.00\quad \textbf{Relev.:}~1.00\quad \textbf{Ctx.Rel.:}~0.25\\

\noindent\textbf{Antwort:} Based on the provided rules context, no, you cannot use Action Surge to cast two leveled spells in one turn. The Spellcasting rule states that you can't cast another spell during the same turn, except for a cantrip with a casting time of 1 action. Additionally, the Sorcerer class rules mention that \ldots

\noindent\textbf{Abgerufene Chunks:}
\begin{itemize}[noitemsep, topsep=2pt]
  \item[\textbf{\#1}] \texttt{[0.0167]} \textit{Spells/# Spellcasting} -- You can't cast another spell during the same turn, except for a cantrip with a c\ldots
  \item[\textbf{\#2}] \texttt{[0.0164]} \textit{Classes/Sorcerer} -- As a bonus action on your turn, you can expend one spell slot and gain a number \ldots
  \item[\textbf{\#3}] \texttt{[0.0161]} \textit{Classes/Sorcerer} -- You can transform unexpended sorcery points into one spell slot as a bonus actio\ldots
  \item[\textbf{\#4}] \texttt{[0.0159]} \textit{Monsters/Lich} -- Only one legendary action option can be used at a time and only at the end of an\ldots
  \item[\textbf{\#5}] \texttt{[0.0156]} \textit{Classes/Sorcerer} -- \#\#\#\# Quickened Spell  When you cast a spell that has a casting time of 1 action,\ldots
\end{itemize}
\vspace{4pt}
\hrule
\vspace{6pt}
